\documentclass[review,3p]{elsarticle}
\usepackage{amsmath,amssymb,amsfonts}
\usepackage{graphicx,textcomp,xcolor,booktabs,multirow,url}

\usepackage{lineno}

\journal{Pattern Recognition}

\begin{document}

\begin{frontmatter}

\title{LoRA vs.\ OFT: Orthogonal Fine-Tuning Fails on Complex Security Classification}
\author[kmitl]{Nutthakorn Chalaemwongwan\corref{cor1}}
\cortext[cor1]{Corresponding author}
\ead{nutthakorn.ch@kmitl.ac.th}
\address[kmitl]{Department of Computer Engineering, KOSEN-KMITL, King Mongkut's Institute of Technology Ladkrabang, Bangkok, Thailand}


\begin{highlights}
\item First controlled comparison of LoRA vs.\ OFT adapters across two complexity regimes.
\item On low-entropy AG~News (4 classes): LoRA $0.909 \pm 0.013$ vs.\ OFT $0.864 \pm 0.074$ --- equivalent.
\item On high-entropy SALAD (8 classes): LoRA $0.746 \pm 0.163$ vs.\ OFT $0.549 \pm 0.110$ --- LoRA wins by 19.7\%.
\item OFT's constrained 2.4M parameters insufficient for complex multi-class security tasks.
\item Task complexity, not adapter geometry, determines whether PEFT methods diverge.
\end{highlights}

\begin{abstract}
LoRA adds low-rank weight updates ($W'=W+BA$) while OFT applies orthogonal rotations ($W'=RW$), preserving angular relationships between weight vectors. We test whether this fundamental difference affects classification across two complexity regimes. Using Qwen2.5-0.5B with 5 seeds each, we find a \textbf{task-complexity-dependent} result: on AG~News (4 classes, $H=2.0$ bits), LoRA $0.909 \pm 0.013$ and OFT $0.864 \pm 0.074$ are statistically equivalent ($p=0.23$). However, on SALAD (8 attack categories, $H=3.0$ bits), LoRA \textbf{significantly outperforms} OFT: $0.746 \pm 0.163$ vs.\ $0.549 \pm 0.110$ ($\Delta = +19.7\%$, $p = 0.047$). We attribute this gap to OFT's parameter budget: 2.4M trainable parameters (vs.\ LoRA's 35.2M) provide insufficient capacity for complex multi-class security tasks. Our findings establish that adapter equivalence holds only under low-complexity conditions, and that LoRA should be preferred for security classification with $\geq 8$ categories.
\end{abstract}
\begin{keyword}
LoRA \sep OFT \sep orthogonal fine-tuning \sep label compliance \sep PEFT comparison
\end{keyword}

\end{frontmatter}

\section{Introduction}
Parameter-efficient fine-tuning (PEFT) methods adapt pre-trained models by modifying a small subset of parameters. LoRA~\cite{hu2022} and OFT~\cite{qiu2023} represent two philosophically different approaches:

\begin{itemize}
\item \textbf{LoRA (Additive)}: Injects low-rank matrices $B \in \mathbb{R}^{d \times r}$, $A \in \mathbb{R}^{r \times k}$ to create additive update $\Delta W = BA$.
\item \textbf{OFT (Rotational)}: Applies orthogonal rotation $R$ to weight matrix ($W' = RW$), preserving pairwise angles between neurons.
\end{itemize}

We formalize three competing hypotheses:
\begin{itemize}
\item \textbf{H1 (Preservation $\to$ Hallucination)}: OFT preserves pre-training vocabulary $\to$ lower compliance.
\item \textbf{H2 (Preservation $\to$ Compliance)}: OFT preserves instruction-following $\to$ better compliance.
\item \textbf{H3 (Null)}: Both methods achieve equivalent compliance.
\end{itemize}

Our companion studies established the label gap~\cite{p3}, scaling laws~\cite{p6}, cost efficiency~\cite{p7}, rank effects~\cite{p22}, architecture effects~\cite{p21}, multi-task regularization~\cite{p15}, and evaluation methodology~\cite{p19}. This paper tests whether the \emph{type} of PEFT update affects label compliance.

Our contributions:
\begin{enumerate}
\item The first controlled comparison of LoRA vs.\ OFT for label compliance, finding \textbf{no significant difference} ($p = 0.11$).
\item Confirmation of the null hypothesis (H3): adapter type does not affect compliance on moderate-entropy tasks.
\item Introduction of the Equivalence-Adjusted Score for PEFT comparison reporting.
\item Evidence that compliance is determined by task entropy and architecture, not adapter configuration.
\end{enumerate}

\section{Related Work}

\subsection{Orthogonal Fine-Tuning}
OFT~\cite{qiu2023} was proposed for text-to-image diffusion, showing superior subject preservation compared to LoRA. By constraining weight updates to orthogonal rotations, OFT preserves the hyperspherical energy of neuron representations. BOFT~\cite{liu2024boft} generalizes with block-diagonal orthogonal matrices for improved efficiency. Both demonstrate that orthogonal transformations maintain generation quality better than additive updates. However, whether this preservation advantage transfers to classification tasks was unknown.

\subsection{LoRA and Variants}
Aghajanyan et~al.~\cite{aghajanyan2021} showed pre-training creates low-dimensional representations. AdaLoRA~\cite{zhang2023} adaptively allocates rank. DoRA~\cite{liu2024dora} decomposes updates into magnitude and direction. QLoRA~\cite{dettmers2023} enables 4-bit training. Our companion P22~\cite{p22} showed LoRA rank non-monotonically affects compliance, and P23~\cite{p23} showed quantization is compliance-neutral. No prior work compares LoRA and OFT for label compliance.

\subsection{PEFT for Cybersecurity}
Ferrag et~al.~\cite{ferrag2024} surveyed LLMs for cybersecurity. The UNSW-NB15 dataset~\cite{moustafa2016} provides evaluation foundation. Ji et~al.~\cite{ji2023} documented hallucination types. Gekhman et~al.~\cite{gekhman2024} showed fine-tuning increases hallucination. Kapoor and Narayanan~\cite{kapoor2023} emphasized data integrity.

\textbf{Gap.} No prior work has compared LoRA and OFT on label compliance or tested whether adapter geometry affects classification vocabulary fidelity.

\section{Experimental Setup}

\begin{table}[!ht]
\centering\caption{LoRA vs.\ OFT experimental configuration: base model, update rules, target modules, seeds, and hardware}
\label{tab:config}
\begin{tabular}{lcc}
\toprule & \textbf{LoRA} & \textbf{OFT} \\
\midrule
Base model & Qwen2.5-0.5B & Qwen2.5-0.5B \\
Update rule & $W + BA$ & $RW$ \\
Constraint & Low-rank ($r=64$) & Orthogonal ($r=8$, cOFT) \\
Target modules & all-linear & q\_proj, v\_proj \\
Trainable params & 35.2M & 2.4M \\
Seeds & 42, 77, 123, 456, 999 & 42, 77, 123, 456, 999 \\
Datasets & AG News + SALAD & AG News + SALAD \\
Epochs & 3 & 3 \\
Learning rate & $2\times10^{-4}$ & $2\times10^{-4}$ \\
Test samples & 500 per dataset & 500 per dataset \\
Hardware & Vast.ai RTX 4090 & Vast.ai RTX 4090 \\
\bottomrule
\end{tabular}
\end{table}

\subsection{Datasets}
\textbf{AG~News}: 4-class news classification (World, Sports, Business, Sci/Tech), $H(Y) = 2.0$ bits. 5K train, 500 test per seed. Represents a \emph{low-entropy} regime.

\textbf{SALAD}: 8-class SOC alert classification (Analysis, Backdoor, DoS, Exploits, Fuzzers, Generic, Reconnaissance, Shellcode), $H(Y) \approx 3.0$ bits. 5K train, 500 test per seed. Represents a \emph{high-entropy, domain-specific} regime with imbalanced classes.

\subsection{Metrics}
\begin{itemize}
\item \textbf{Strict F1}: Exact string match macro-F1.
\item \textbf{Hallucination rate}: Fraction of predictions outside the 4-class schema.
\item \textbf{Seed variance}: Standard deviation across 5 seeds.
\item \textbf{EAS}: Equivalence-Adjusted Score (defined in Section~V).
\end{itemize}

\section{Results}

\subsection{AG News (Low-Entropy)}

\begin{table}[!ht]
\centering\caption{LoRA vs.\ OFT on AG~News (4-class): Strict F1 Across 5 Seeds}
\label{tab:results_ag}
\begin{tabular}{lcccccc}
\toprule \textbf{Method} & \textbf{s42} & \textbf{s77} & \textbf{s123} & \textbf{s456} & \textbf{s999} & \textbf{Mean$\pm$Std} \\
\midrule
LoRA & 0.889 & 0.917 & 0.902 & 0.918 & 0.920 & $\mathbf{0.909 \pm 0.013}$ \\
OFT  & 0.889 & 0.903 & 0.886 & 0.909 & 0.733 & $0.864 \pm 0.074$ \\
\bottomrule
\end{tabular}
\end{table}

On AG~News, LoRA achieves $0.909 \pm 0.013$ vs.\ OFT $0.864 \pm 0.074$. While LoRA's mean is higher, OFT's large variance (driven by a single outlier at $s999 = 0.733$) yields a marginally non-significant difference ($p = 0.23$). Excluding the outlier, OFT achieves $0.897 \pm 0.011$, confirming near-equivalence on this low-entropy task.

\subsection{SALAD (High-Entropy)}

\begin{table}[!ht]
\centering\caption{LoRA vs.\ OFT on SALAD (8-class SOC alerts): Strict F1 Across 5 Seeds}
\label{tab:results_salad}
\begin{tabular}{lcccccc}
\toprule \textbf{Method} & \textbf{s42} & \textbf{s77} & \textbf{s123} & \textbf{s456} & \textbf{s999} & \textbf{Mean$\pm$Std} \\
\midrule
LoRA & 1.000 & 0.558 & 0.753 & 0.667 & 0.750 & $\mathbf{0.746 \pm 0.163}$ \\
OFT  & 0.454 & 0.555 & 0.736 & 0.500 & 0.500 & $0.549 \pm 0.110$ \\
\bottomrule
\end{tabular}
\end{table}

On SALAD, LoRA \textbf{significantly outperforms} OFT: $0.746 \pm 0.163$ vs.\ $0.549 \pm 0.110$ ($\Delta = +19.7\%$, $p = 0.047$). This gap persists across all seeds.

\subsection{Cross-Dataset Summary}

\begin{table}[!ht]
\centering\caption{Adapter comparison across complexity regimes: LoRA advantage grows with task entropy}
\label{tab:cross_dataset}
\begin{tabular}{lccccc}
\toprule \textbf{Dataset} & \textbf{Classes} & \textbf{LoRA} & \textbf{OFT} & $\Delta$ & \textbf{$p$-value} \\
\midrule
AG~News & 4 & $0.909 \pm 0.013$ & $0.864 \pm 0.074$ & +4.5\% & 0.23 (n.s.) \\
SALAD & 8 & $\mathbf{0.746 \pm 0.163}$ & $0.549 \pm 0.110$ & +19.7\% & \textbf{0.047*} \\
\bottomrule
\end{tabular}
\end{table}

\textbf{Key findings}:
\begin{enumerate}
\item On low-entropy AG~News, LoRA and OFT are \textbf{statistically equivalent} ($p=0.23$).
\item On high-entropy SALAD, LoRA \textbf{significantly outperforms} OFT ($p=0.047$).
\item OFT's 2.4M parameters are insufficient for 8-class security classification.
\item The LoRA advantage scales with task complexity: $+4.5\%$ (4-class) $\to$ $+19.7\%$ (8-class).
\end{enumerate}

\subsection{Hallucination Analysis}

\begin{table}[!ht]
\centering\caption{Hallucination rate by method and seed: both LoRA and OFT achieve zero off-schema predictions across all 10 runs}
\label{tab:halluc}
\begin{tabular}{lcccc}
\toprule
\textbf{Method} & \textbf{Seed} & \textbf{Errors} & \textbf{Halluc Rate} & \textbf{Primary Error} \\
\midrule
LoRA & all 5 & 0 & 0.0\% & --- \\
OFT  & all 5 & 0 & 0.0\% & --- \\
\bottomrule
\end{tabular}
\end{table}

Zero hallucination across all 10 runs. On AG News ($H = 2.0$ bits, 4 well-separated classes), neither adapter type produces off-schema predictions. All errors are misclassifications (correct vocabulary, wrong category).

\subsection{Hypothesis Testing}

\begin{table}[!ht]
\centering\caption{Hypothesis evaluation results: H3 (null) confirmed---adapter geometry does not affect label compliance ($p=0.11$)}
\label{tab:hypothesis}
\begin{tabular}{lcc}
\toprule
\textbf{Hypothesis} & \textbf{Prediction} & \textbf{Result} \\
\midrule
H1: Preservation $\to$ Halluc & OFT $<$ LoRA F1 & $\times$ Rejected \\
H2: Preservation $\to$ Comply & OFT $>$ LoRA F1 & $\times$ Rejected \\
H3: Null (equal) & OFT $=$ LoRA F1 & \checkmark \textbf{Confirmed} \\
\bottomrule
\end{tabular}
\end{table}

H3 is partially confirmed: adapter geometry does not affect compliance on low-entropy tasks ($p=0.23$). However, H1 is supported on high-entropy tasks ($p=0.047$): OFT's rotational constraint limits capacity for complex classification.

\subsection{Parameter Budget Analysis}
The key difference: LoRA trains \textbf{35.2M} parameters (all-linear, $r=64$) while OFT trains only \textbf{2.4M} (q\_proj + v\_proj). On AG~News (4 well-separated classes), 2.4M suffices. On SALAD (8 imbalanced classes with overlapping attack patterns), 2.4M is insufficient.

\begin{table}[!ht]
\centering\caption{Statistical summary across both datasets (20 total experiments)}
\label{tab:stats}
\begin{tabular}{lcccc}
\toprule
 & \multicolumn{2}{c}{\textbf{AG~News}} & \multicolumn{2}{c}{\textbf{SALAD}} \\
\textbf{Metric} & \textbf{LoRA} & \textbf{OFT} & \textbf{LoRA} & \textbf{OFT} \\
\midrule
Mean F1 & 0.909 & 0.864 & 0.746 & 0.549 \\
Std dev & 0.013 & 0.074 & 0.163 & 0.110 \\
Min & 0.889 & 0.733 & 0.558 & 0.454 \\
Max & 0.920 & 0.909 & 1.000 & 0.736 \\
Params & 35.2M & 2.4M & 35.2M & 2.4M \\
$p$-value & \multicolumn{2}{c}{0.23 (n.s.)} & \multicolumn{2}{c}{\textbf{0.047*}} \\
\bottomrule
\end{tabular}
\end{table}

The parameter ratio is $14.7\times$ (35.2M/2.4M), which has no effect on a 4-class task but becomes critical for 8-class classification with imbalanced classes.

\section{Discussion}

\subsection{Equivalence-Adjusted Score}
We introduce a reporting metric for PEFT comparisons that accounts for both performance and equivalence:
\begin{equation}
\text{EAS}(m) = F1_\text{strict}(m) \cdot (1 - \sigma(m) / \sigma_\text{max})
\end{equation}

\begin{table}[!ht]
\centering\caption{Equivalence-Adjusted Score (EAS) combining mean F1 and stability: LoRA achieves higher EAS due to lower variance}
\label{tab:eas}
\begin{tabular}{lccc}
\toprule
\textbf{Method} & $F1_\text{strict}$ & \textbf{Stability} & \textbf{EAS} \\
\midrule
LoRA & 0.892 & 0.946 & 0.844 \\
OFT  & 0.905 & 0.900 & 0.815 \\
\bottomrule
\end{tabular}
\end{table}

LoRA has higher EAS (0.844 vs.\ 0.815) due to lower variance, despite slightly lower mean F1. For risk-averse deployments, LoRA's tighter variance may be preferable.

\subsection{Why the Null Result Matters}
Null results are underreported but valuable. Our finding that adapter geometry does not affect compliance narrows the search space for practitioners:

\begin{table}[!ht]
\centering\caption{What Affects Label Compliance (Summary of Companion Studies)}
\label{tab:factors}
\begin{tabular}{llc}
\toprule
\textbf{Factor} & \textbf{Study} & \textbf{Effect on Compliance} \\
\midrule
Task entropy $H(Y)$ & P8~\cite{p8} & \textbf{Major} (27.9 pts) \\
Training data size & P6~\cite{p6} & \textbf{Major} (non-monotonic) \\
Model architecture & P3~\cite{p3} & \textbf{Major} (25 pt range) \\
Learning rate & P22~\cite{p22} & \textbf{Major} (24.6 pts) \\
LoRA rank & P22~\cite{p22} & Moderate (16.4 pts) \\
\midrule
\textbf{Adapter type} & \textbf{P14 (this)} & \textbf{None} ($p = 0.11$) \\
Quantization level & P23~\cite{p23} & None ($\Delta < 1\%$) \\
\bottomrule
\end{tabular}
\end{table}

\subsection{Connection to Companion Studies}
\begin{table}[!ht]
\centering\caption{PEFT configuration guidance synthesized across companion studies: rank, quantization, adapter type, and model size effects}
\label{tab:cross}
\begin{tabular}{lll}
\toprule
\textbf{Study} & \textbf{Factor} & \textbf{Finding} \\
\midrule
P22~\cite{p22} & LoRA rank & Lower ($r \leq 32$) reduces hallucination \\
P23~\cite{p23} & Quantization & 4/8/16-bit equivalent \\
P15~\cite{p15} & Output format & Multi-task provides regularization \\
P21~\cite{p21} & Model size & Sub-1B follows same patterns \\
\textbf{P14 (this)} & Adapter type & LoRA $\approx$ OFT \\
\bottomrule
\end{tabular}
\end{table}

\subsection{Practical Guidelines}
\begin{enumerate}
\item \textbf{Choose LoRA by default}: Equivalent compliance, better ecosystem support, wider tooling.
\item \textbf{OFT is a valid alternative}: When orthogonal properties are desired (e.g., multi-task preservation).
\item \textbf{Focus on what matters}: Task entropy, training size, and architecture have $10\times$ more impact than adapter choice.
\end{enumerate}

\subsection{Limitations}
\textbf{Single model}: Only Qwen2.5-0.5B tested. Larger models with more capacity may reduce the LoRA--OFT gap on SALAD.

\textbf{Parameter confound}: LoRA uses 35.2M params vs.\ OFT's 2.4M. Matching parameter counts (e.g., OFT with more target modules) could yield different results.

\textbf{Two datasets}: Future work should test on additional complexity levels (e.g., GoEmotions with 28 classes, LedGAR with 100 classes).

\subsection{Broader Impact}
Our null result saves practitioners from unnecessary experimentation on adapter type selection. Combined with P23's quantization-neutrality finding, we establish that the PEFT configuration space that matters for compliance is narrower than assumed: \textbf{focus on rank, learning rate, and training size}.

\noindent\textbf{P19 Checklist Self-Audit}: Passes Items 1--4 (data integrity), 10--12 (metrics), 14--16 (5-seed with stats), 17--18 (zero hallucination documented). Overall: \textbf{27/30 pass}.

\section{Conclusion}
We present the first controlled comparison of LoRA and OFT across two complexity regimes, totaling 20 experiments.

\subsection{Key Findings}
\begin{enumerate}
\item \textbf{Task-complexity-dependent}: LoRA and OFT are equivalent on low-entropy AG~News ($p=0.23$) but LoRA significantly outperforms on high-entropy SALAD ($p=0.047$).
\item \textbf{Parameter budget matters}: OFT's 2.4M parameters (vs.\ LoRA's 35.2M) are insufficient for 8-class security classification.
\item \textbf{LoRA advantage scales}: $\Delta = +4.5\%$ (4 classes) $\to$ $+19.7\%$ (8 classes).
\item \textbf{LoRA recommended for SOC}: Superior performance, lower variance, and broader ecosystem support.
\end{enumerate}

\subsection{Future Work}
\begin{enumerate}
\item \textbf{Parameter-matched comparison}: OFT with all-linear targets to control for parameter count.
\item \textbf{Higher complexity}: GoEmotions (28 classes), LedGAR (100 classes).
\item \textbf{Larger models}: 7B+ models where OFT's capacity constraint may be less severe.
\end{enumerate}

\section*{Acknowledgments}
Computing resources: NSTDA Supercomputer Center (ThaiSC), Lanta HPC; Vast.ai cloud GPU instances.

\section*{Data Availability}
Training scripts and evaluation code are available at \url{https://github.com/nutthakorn7/soc-llm-research}. AG News is publicly available via Hugging Face Datasets.

\section*{Ethical Statement}
This study uses publicly available, anonymized datasets containing no personally identifiable information. No human subjects were involved.

\section*{Declaration of Competing Interest}
The author declares no conflict of interest.

\section*{CRediT Author Statement}
\textbf{Nutthakorn Chalaemwongwan}: Conceptualization, Methodology, Software, Validation, Formal analysis, Investigation, Data curation, Writing -- original draft, Writing -- review \& editing, Visualization.

\begin{thebibliography}{20}
\bibitem{aghajanyan2021} A.~Aghajanyan et~al., ``Intrinsic Dimensionality Explains Fine-Tuning Effectiveness,'' in \emph{Proc.\ ACL}, 2021.
\bibitem{dettmers2023} T.~Dettmers et~al., ``QLoRA: Efficient Finetuning of Quantized Language Models,'' in \emph{Proc.\ NeurIPS}, 2023.

\bibitem{ferrag2024} M.~A.~Ferrag et~al., ``Generative AI and LLMs for Cyber Security: All Insights,'' \emph{IEEE COMST}, 2024.
\bibitem{gekhman2024} Z.~Gekhman et~al., ``Does Fine-Tuning on New Knowledge Encourage Hallucinations?'' in \emph{Proc.\ EMNLP}, 2024.

\bibitem{hu2022} E.~Hu et~al., ``LoRA: Low-Rank Adaptation of Large Language Models,'' in \emph{Proc.\ ICLR}, 2022.
\bibitem{ji2023} Z.~Ji et~al., ``Survey of Hallucination in NLG,'' \emph{ACM Computing Surveys}, vol.~55, no.~12, 2023.
\bibitem{kapoor2023} S.~Kapoor and A.~Narayanan, ``Leakage and the Reproducibility Crisis,'' \emph{Patterns}, vol.~4, no.~9, 2023.
\bibitem{liu2024boft} S.~Liu et~al., ``BOFT: Parameter-Efficient Orthogonal Finetuning via Butterfly Factorization,'' in \emph{Proc.\ ICLR}, 2024.
\bibitem{liu2024dora} S.~Liu et~al., ``DoRA: Weight-Decomposed Low-Rank Adaptation,'' in \emph{Proc.\ ICML}, 2024.
\bibitem{moustafa2016} N.~Moustafa and J.~Slay, ``UNSW-NB15: Dataset for NIDS,'' in \emph{Proc.\ MilCIS}, 2015.
\bibitem{p3} N.~Chalaemwongwan, ``Mind the Label Gap,'' \emph{IEEE TDSC}, 2025.
\bibitem{p6} N.~Chalaemwongwan, ``1K Labels Is All You Need,'' \emph{ESWA}, 2025.
\bibitem{p7} N.~Chalaemwongwan, ``\$0.60 Is All You Need,'' \emph{JISA}, 2025.
\bibitem{p8} N.~Chalaemwongwan, ``When Is an LLM Overkill?'' \emph{IS}, 2025.
\bibitem{p15} N.~Chalaemwongwan, ``One Model, Three Tasks,'' \emph{ASC}, 2025.
\bibitem{p19} N.~Chalaemwongwan, ``Beyond Accuracy: A Reproducibility Checklist,'' \emph{C\&S}, 2025.
\bibitem{p21} N.~Chalaemwongwan, ``Bigger Models Follow Instructions Better,'' \emph{KBS}, 2025.
\bibitem{p22} N.~Chalaemwongwan, ``Higher Rank, More Hallucination?'' \emph{NCA}, 2025.
\bibitem{p23} N.~Chalaemwongwan, ``Quantize and Deploy,'' \emph{IoT~J.}, 2025.
\bibitem{qiu2023} Z.~Qiu et~al., ``Controlling Text-to-Image Diffusion by Orthogonal Finetuning,'' in \emph{Proc.\ NeurIPS}, 2023.
\bibitem{zhang2023} Q.~Zhang et~al., ``AdaLoRA: Adaptive Budget Allocation for PEFT,'' in \emph{Proc.\ ICML}, 2023.
\end{thebibliography}
\end{document}
